% Paper 2 (ReSIDS v2) -- datasets and domain profiles. Impersonal style.
% Citation keys are in docs/refs_v2.bib. Conditions of use come from the audit
% (Section~\ref{sec:dataset-audit}; reports in results/*_audit.txt).
% DECISIONS encoded below (confirm): CICIDS2017 'Infiltration - Portscan' merged into
% PortScan; 'Attempted' flows labelled benign; Heartbleed excluded; Init Win Bytes
% excluded.
% 2026-10-07: CICIoT2023 dropped (labels mark the capture session in four of seven
% categories; audit). ToN_IoT and Morris-2/3/4 rejected. Electra Modbus retained as the
% third domain with four content-observable categories (scripts/cross_domain_electra.py:
% GL k>=2 recall 99.94%, 0 false positives on 4.17 M normal packets; FL 51.7% with 4 nodes).
% IoT-23 still under audit. The count-based tick and attribution "none" were exercised
% only by CICIoT2023.

\subsection{Datasets and domain profiles}
\label{sec:datasets}

ReSIDS~v2 claims generality for its core---the event model, the authority rule, the two
Disaster-FD cadences and the reconciliation contract---and not for its detectors.
Detectors trained on one intrusion dataset do not carry over to
another~\cite{dhooge2020interdataset}. In every domain the specialists are retrained,
and only the profile (Section~\ref{sec:api-events}) changes. The evaluation therefore
needs domains that differ in what the core abstracts away, rather than datasets that
maximise accuracy. Three criteria guided the choice:
(i)~the \emph{domain}, i.e., the network the agent protects and the shape of its traffic;
(ii)~the \emph{attributes the core consumes besides the features}, namely whether the
traffic carries time and whether an alarm can be attributed to an emitter; and
(iii)~\emph{verified data integrity}. Label inaccuracy and the base-rate fallacy are
among the most frequent pitfalls of learning-based security
evaluation~\cite{axelsson2000baserate,sommer2010outside,arp2022dos}, so every candidate
was audited before use (Section~\ref{sec:dataset-audit}). Three public datasets passed:
an IEC~61850 substation, an enterprise network and the Modbus control network of a
railway traction substation. Table~\ref{tab:datasets} summarises them.

\begin{table}[t]
\centering
\caption{Datasets, the profile each one binds, and its role in the evaluation.}
\label{tab:datasets}
\small
\begin{tabular}{@{}llllll@{}}
\toprule
Dataset & Domain & Local tick & Attribution & Label set & Role \\
\midrule
ERENO~\cite{quincozes2023ereno} & IEC 61850 substation & time & synthetic only$^{a}$ &
  \texttt{ereno-7} & case study \\
CICIDS2017-C~\cite{engelen2021troubleshooting} & enterprise network & time &
  source address & \texttt{cicids2017c-7} & attribution and time \\
Electra Modbus~\cite{perales2019electra} & industrial control network & time &
  source MAC address & \texttt{electra-4} & industrial traffic \\
\bottomrule
\multicolumn{6}{@{}l}{\footnotesize $^{a}$ All messages carry the legitimate publisher's
identity; see the ERENO paragraph.}
\end{tabular}
\end{table}

Every domain uses two specialists per category of its label set: $N=14$ for ERENO and
the corrected CICIDS2017, $N=8$ for Electra. Identifiers (addresses, ports, timestamps,
flow and message identities) feed the profile, i.e., the tick and the attribution, and
never the features. This avoids the shortcuts that identifiers create in benchmark
datasets~\cite{flood2024smells}.

\paragraph{ERENO (case study).}
ERENO~\cite{quincozes2023ereno} generates IEC~61850 GOOSE and Sampled Values traffic for
digital substations, with seven attacks against the publish--subscribe messages that
carry protection signals. It is the case study inherited from v1, the only domain whose
traffic is time-critical, and the most demanding of the three. The authors' train/test
split is kept: each periodic GOOSE frame accompanies thousands of SV frames, so a random
re-split would leak. Publisher identities (MAC addresses, \texttt{gocbRef},
\texttt{goID}) and the three absolute-time markers are removed from the features, so
that a specialist learns how an attack behaves and not who performs it or when. The
audit sets two limits. First, the attacks impersonate the single legitimate publisher,
so source attribution would point at the victim; intrusion-driven isolation is
therefore exercised only in synthetic scenarios. Second, the benign volume available at
window level, about twelve minutes, bounds the precision that can be certified at low
prevalence (Section~\ref{sec:api-fp}).

\paragraph{CICIDS2017, corrected (attribution and time).}
CICIDS2017~\cite{sharafaldin2018toward} is the most cited network intrusion benchmark,
but its original labels and flow features contain errors that change detection
results~\cite{engelen2021troubleshooting,lanvin2023errors,liu2022error}. The corrected
re-extraction of Engelen et al.~\cite{engelen2021troubleshooting,liu2022error} is used,
with its five days (2{,}099{,}976 flows). It is the only domain that provides, on
captured traffic, three properties the core needs. Real flow timestamps let the
time-based local tick run on live traffic. Source addresses point at the emitter, either
the attacker's gateway or a compromised internal host, so the \texttt{device\_address}
attribution, and therefore intrusion-driven isolation, can be exercised. A benign-only
day (Monday, eight hours) extends the benign volume available for the prevalence
analysis well beyond that of ERENO. The label set \texttt{cicids2017c-7} groups the
classes into DoS, DDoS, port scan, FTP/SSH brute force, web attack, botnet and
infiltration; Heartbleed (11 flows) is excluded. The audit sets three conditions of use.
Flows labelled ``Infiltration -- Portscan'' are merged into port scan, because 91{,}253
of the scan flows are identical under both labels. The initial TCP window sizes are
excluded from the features, because they fingerprint the scanning host. And the 11{,}979
flows marked \emph{Attempted}, whose attacks delivered no payload, are labelled benign.
The web and infiltration categories keep only 104 and 36 flows, and their figures are
reported with confidence intervals.

\paragraph{Electra Modbus (industrial traffic).}
Electra~\cite{perales2019electra} records the control network of an electric traction
substation built after one in service on a high-speed railway: five programmable logic
controllers and a SCADA system that communicate over Modbus/TCP and S7Comm. The Modbus
release is used, with 16{,}289{,}277 packets over about fourteen hours. It adds the
domain the other two lack, industrial control traffic with request--response polling, and
it is the dataset closest to operational prevalence: the attacks retained below make up
0.41\% of the packets once the excluded types are removed. Packet timestamps drive the time-based tick, and every attack packet has the
man-in-the-middle host as its source or destination, identified by its MAC address, so
\texttt{device\_address} attribution is exercised in a second domain. The audit sets one condition of use. Three of the seven
attack labels mark the path of a packet rather than its content: packets relayed
unchanged by the man-in-the-middle host (\texttt{MITM\_UNALTERED}), replayed packets and
most read requests issued by that host have the same content as normal packets and are
told apart only by the host's MAC address. Those three types are excluded. The label set
\texttt{electra-4} keeps the four types visible in the packet content: reconnaissance,
unauthorised writes, forced errors and modified responses. Features are the content fields
of each packet (request flag, function code, error code, register address and value);
addresses and time feed the profile only. The forced-error category keeps 1{,}129 packets,
and its figures are reported with confidence intervals.

\paragraph{Datasets not used.}
The original CICIDS2017 release is not used, because of the label errors documented
above; neither is CSE-CIC-IDS2018, which shares its generation methodology and its error
classes~\cite{liu2022error}. NSL-KDD~\cite{tavallaee2009nslkdd}, still common in recent
work, is not considered: its records carry neither timestamps nor addresses, and it
inherits the artefacts of the 1998 DARPA traffic~\cite{mchugh2000darpa}. CICIoT2023~\cite{neto2023ciciot2023}, first selected as the
IoT domain, was rejected: in four of its seven categories the labels mark the capture
session rather than the attack, and specialists trained on them flag benign traffic as
attack. ToN\_IoT~\cite{moustafa2021toniot}, audited to replace it, has consistent labels
but a narrow normal class: under the same protocol, 18\% of its normal flows are flagged.
Re-extractions of the same captures inherit these defects and are not considered either:
CIC-ToN-IoT~\cite{cic2021toniot} re-extracts the ToN\_IoT captures with CICFlowMeter, and
NF-UQ-NIDS~\cite{sarhan2021netflow} merges NetFlow versions of four releases, ToN\_IoT and
CSE-CIC-IDS2018 among them, so that the network a flow comes from also identifies its
source dataset.
Thirteen further candidates, from industrial, medical, aerial, power-system, gas-pipeline,
water and industrial-control settings, failed the audit or added no property the three
retained datasets lack (Section~\ref{sec:dataset-audit}). One of them, HAI, has sound
labels and genuine process data but is built for one-class anomaly detection, as are SWaT
and WADI~\cite{mathur2016swat,ahmed2017wadi}, available on request; using them is left to
a one-class variant of the specialists.

\paragraph{Threats to validity.}
The corrected CICIDS2017 is a laboratory capture, Electra a testbed reproducing an
industrial installation, and ERENO synthetic traffic generated by deterministic attack
rules. The four Electra categories are separable by packet content, so that domain tests
the architecture under node loss rather than the difficulty of detection. A detector evaluated on them may not reflect
operational conditions~\cite{sommer2010outside,arp2022dos}. Categories with fewer than
about a hundred positive flows yield per-category figures with wide confidence
intervals; they are reported, but no conclusion is drawn from them. Results are not
compared across datasets: the claim under test is that the same core runs unchanged
under different profiles, not that a model trained in one domain detects attacks in another.
